%! Expressions and Operations — Unit 1, Lesson 1.1 — Slides (Beamer)
\documentclass[aspectratio=169, 11pt]{beamer}
\usepackage{algebra2-beamer}   % provides \forestheader{} and \sectionlabel[color]{}

\begin{document}

% TITLE SLIDE (CourseName is not defined in beamer, so the name is literal)
{
\setbeamercolor{background canvas}{bg=forest}
\begin{frame}[plain]
  \vspace{2.8cm}\hspace{0.55cm}%
  \begin{minipage}{0.9\textwidth}
    {\color{goldacc}\bfseries\small LESSON 1.1}\\[0.5cm]
    {\color{white}\bfseries\LARGE Expressions and Operations}\\[1.2cm]
    {\color{forestmid}\small Algebra 2: Shepherd~$\cdot$~Unit 1: Foundations --- Algebra 1 Review}
  \end{minipage}
\end{frame}
}

% HOOK
\begin{frame}[t, plain]
  \forestheader{Equivalent, or not? Ninety seconds. No rules.}
  \sectionlabel{Hook}
  \begin{columns}[T]
    \begin{column}{0.44\textwidth}
      \centering
      \renewcommand{\arraystretch}{1.6}
      \begin{tabular}{c l c l}
        (a) & $2^3\cdot2^4$ & vs. & $2^{12}$ \\
        (b) & $(x+3)(x-3)$ & vs. & $x^2-9$ \\
        (c) & $(x+5)^2$ & vs. & $x^2+25$ \\
        (d) & $\sqrt{9+16}$ & vs. & $\sqrt9+\sqrt{16}$ \\
      \end{tabular}
    \end{column}
    \begin{column}{0.52\textwidth}
      \begin{itemize}\setlength{\itemsep}{6pt}
        \item Three of these die to a \textbf{single number}.
        \item (a): $2^3\cdot2^4=128$, not $4096$.
        \item (c): at $x=1$, \; $36 \ne 26$.
        \item (d): $\sqrt{25}=5$, but $3+4=7$.
        \item (b) survives \emph{every} value you try.
      \end{itemize}
    \end{column}
  \end{columns}
  \begin{block}{The asymmetry --- keep it all year}
    One bad value \textbf{disproves} equivalence forever. No number of good values ever
    \textbf{proves} it. For proof you need algebra: multiply $(x+3)(x-3)$ out and you have settled
    every value at once.
  \end{block}
\end{frame}

% SECTION 1 — translating and evaluating
\begin{frame}[t, plain]
  \forestheader{An expression is a recipe, not a question}
  \sectionlabel{1. Translate \& evaluate}
  \begin{columns}[T]
    \begin{column}{0.50\textwidth}
      \renewcommand{\arraystretch}{1.35}
      \begin{tabular}{@{}l l@{}}
        Seven less than twice a number & $2n-7$ \\
        Quotient of a number and 4, & \\
        \quad increased by 9 & $\tfrac n4+9$ \\
        Five times the quantity & \\
        \quad of a number decreased by 2 & $5(n-2)$ \\
      \end{tabular}
      \vspace{8pt}

      \textbf{Two traps:} ``less than'' \textbf{reverses} the order; ``the quantity'' means
      \textbf{parentheses}.
    \end{column}
    \begin{column}{0.46\textwidth}
      \begin{itemize}\setlength{\itemsep}{5pt}
        \item An expression is never \emph{solved} --- it is \textbf{evaluated}, once you supply
          values.
        \item $|\,\cdot\,|$ and $\sqrt{\phantom{x}}$ are \textbf{grouping symbols}: finish the
          inside first.
        \item At $x=-3$: \quad
          $(-x)^2 = {\color{forest}\mathbf{9}}$ \quad but \quad $-x^2 = {\color{goldacc}\mathbf{-9}}$
      \end{itemize}
    \end{column}
  \end{columns}
  \begin{block}{The most common evaluation error in the course}
    An exponent binds \textbf{only to what it touches}. Parentheses are what change that --- and
    this one error comes back in the quadratic formula, in end behavior, and in every graphing
    lesson.
  \end{block}
\end{frame}

% SECTION 2 — laws of exponents
\begin{frame}[t, plain]
  \forestheader{The exponent laws are counting, not memorizing}
  \sectionlabel{2. Derived, not memorized}
  \begin{columns}[T]
    \begin{column}{0.50\textwidth}
      \renewcommand{\arraystretch}{1.35}
      \begin{tabular}{@{}l l@{}}
        $a^3a^2$ & five $a$'s \; $\Rightarrow$ \; exponents \textbf{add} \\
        $\dfrac{a^5}{a^2}$ & two cancel \; $\Rightarrow$ \; exponents \textbf{subtract} \\
        $(a^3)^2$ & a pile, twice \; $\Rightarrow$ \; exponents \textbf{multiply} \\
        $(ab)^n$ & every factor inside gets it \\
      \end{tabular}
      \vspace{8pt}

      Write the factors out. Count them. That is the whole derivation.
    \end{column}
    \begin{column}{0.46\textwidth}
      \begin{block}{Forced, not invented}
        $\dfrac{a^3}{a^3}$ is both $a^0$ and $1$ \; $\Rightarrow$ \; $a^0=1$ \\[4pt]
        $\dfrac{a^2}{a^5}$ is both $a^{-3}$ and $\dfrac{1}{a^3}$ \; $\Rightarrow$ \;
        $a^{-n}=\dfrac{1}{a^n}$
      \end{block}
      \begin{itemize}\setlength{\itemsep}{4pt}
        \item $a^{-2}$ is a \textbf{reciprocal}, not a negative number.
        \item Final answers carry \textbf{positive} exponents only.
      \end{itemize}
    \end{column}
  \end{columns}
  \vspace{2pt}
  \centering
  $3x^2\cdot4x^3=12x^5$ --- coefficients \textbf{multiply}, exponents \textbf{add}. \quad
  $\dfrac{12x^5y^{-2}}{3x^2y^4}=\dfrac{4x^3}{y^6}$
\end{frame}

% SECTION 3 — polynomials and factoring
\begin{frame}[t, plain]
  \forestheader{Factoring is multiplying, walked backwards}
  \sectionlabel{3. Polynomials \& factoring}
  \begin{columns}[T]
    \begin{column}{0.48\textwidth}
      \begin{itemize}\setlength{\itemsep}{5pt}
        \item Subtracting? Distribute the minus to \textbf{every} term --- not just the first.
        \item $a^2-b^2=(a+b)(a-b)$
        \item $(a+b)^2=a^2+2ab+b^2$ \\
          {\small Hook (c) settled: the missing $10x$ was the whole point.}
        \item Every factorization is \textbf{checkable by expanding}. So check it.
      \end{itemize}
    \end{column}
    \begin{column}{0.48\textwidth}
      \begin{block}{The order, every single time}
        \textbf{(1)} GCF \emph{first}, always. \\
        \textbf{(2)} Then the pattern or the trinomial method. \\
        \textbf{(3)} Then ask: does any factor \emph{still} factor?
      \end{block}
      \centering
      $2x^2-18 \;\to\; 2(x^2-9) \;\to\; {\color{forest}\mathbf{2(x+3)(x-3)}}$

      \vspace{4pt}
      {\small Stopping at $2(x^2-9)$ is factored --- but not \textbf{completely}.}
    \end{column}
  \end{columns}
  \begin{block}{Why ``completely'' is not fussiness}
    In Lesson 1.2 you will factor in order to \textbf{solve}. Stop early there and you do not lose a
    step --- you lose a \textbf{solution}.
  \end{block}
\end{frame}

% SECTION 4 — radicals and rational exponents
\begin{frame}[t, plain]
  \forestheader{A radical and a fractional exponent are one idea, spelled twice}
  \sectionlabel{4. Radicals \& rational exponents}
  \begin{columns}[T]
    \begin{column}{0.50\textwidth}
      \renewcommand{\arraystretch}{1.4}
      \begin{tabular}{@{}l l l@{}}
        $\sqrt{72}$ & $=\sqrt{36\cdot2}$ & $=6\sqrt2$ \\
        $\sqrt[3]{54}$ & $=\sqrt[3]{27\cdot2}$ & $=3\sqrt[3]{2}$ \\
      \end{tabular}
      \vspace{6pt}

      Same method. Only the \textbf{index} changed.
      \vspace{8pt}

      \begin{itemize}\setlength{\itemsep}{4pt}
        \item $3\sqrt5+2\sqrt5=5\sqrt5$ \quad --- like radicals.
        \item $3\sqrt5+2\sqrt3$ \quad --- will not combine. Ever.
        \item $\sqrt{12}+\sqrt{27}=2\sqrt3+3\sqrt3=5\sqrt3$
      \end{itemize}
    \end{column}
    \begin{column}{0.46\textwidth}
      \begin{block}{Two spellings}
        $a^{1/2}=\sqrt{a}$ \qquad $a^{1/3}=\sqrt[3]{a}$
      \end{block}
      \centering
      $\sqrt[3]{8}=8^{1/3}=(2^3)^{1/3}=2^{3\cdot\frac13}=2$

      \vspace{8pt}
      \raggedright
      The \textbf{power-of-a-power law} did that. The exponent spelling is what lets the exponent
      laws reach radicals at all --- and it is the notation all of \textbf{Unit 6} is built on.
    \end{column}
  \end{columns}
\end{frame}

% CLOSE
\begin{frame}[t, plain]
  \forestheader{What today actually was}
  \sectionlabel{Closing}
  \begin{columns}[T]
    \begin{column}{0.48\textwidth}
      \begin{block}{Errors that cost the most this year}
        $(x+5)^2 = x^2+25$ \\
        $\sqrt{a+b} = \sqrt a + \sqrt b$ \\[4pt]
        {\small Same mistake, two costumes: assuming an operation distributes when it does not.}
      \end{block}
    \end{column}
    \begin{column}{0.48\textwidth}
      \begin{block}{Habits to leave with}
        Check a factorization by \textbf{expanding} it. \\
        Kill a false claim with \textbf{one value}. \\
        Prove a true one with \textbf{algebra}.
      \end{block}
    \end{column}
  \end{columns}
  \vspace{6pt}
  \begin{itemize}\setlength{\itemsep}{5pt}
    \item \textbf{Exit ticket:} one evaluation, one factor-completely with a check, one SOL-style
      exponent item.
    \item \textbf{Next --- Lesson 1.2, Equations and Inequalities:} every rewriting rule from today
      goes to work. Distributing clears parentheses, factoring plus the zero product property
      \emph{solves} quadratics, and the quadratic formula hands you a radical you can now simplify.
  \end{itemize}
\end{frame}

\end{document}
